\section{Appendix: audit materials}
\label{app:audit-materials}

This appendix collects three audit aids: a dependency map, a denial-condition matrix, and a compact
reading path. It adds no premises to the theorem ladder.

\subsection{Dependency map for the theorem ladder}

\begin{figure}[htbp]
\centering
\begin{tikzpicture}[
  node distance=10mm and 10mm,
  >=Latex,
  every node/.style={font=\small},
  box/.style={draw=accent, rounded corners=2pt, fill=softaccent, align=left, inner sep=6pt, text width=0.29\textwidth},
  line/.style={-Latex, semithick, draw=accent}
]
\node[box] (substrate) {\textbf{Fixed downstream substrate}\\
Definition~\ref{def:fixed-downstream-substrate} and Proposition~\ref{prop:substrate-consequences}.};
\node[box, right=of substrate] (semantic) {\textbf{Measurement-reference quotient}\\
Definitions~\ref{def:admissible-reference-assignment}--\ref{def:mr-presentation};
Theorem~\ref{thm:measurement-reference-quotient}.};
\node[box, below=of substrate] (everett) {\textbf{Everettian non-selection}\\
Corollary~\ref{cor:everettian-nonselection} and Theorem~\ref{thm:everett_nonselection}.};
\node[box, below=of semantic] (confirmation) {\textbf{Confirmation criterion}\\
Theorems~\ref{thm:theory-level-confirmation-criterion}, \ref{thm:fixation-commutation-necessity}, and \ref{thm:reference-stabilization-necessity}.};
\node[box, below=of confirmation] (toy) {\textbf{Toy witness}\\
Theorem~\ref{thm:toy-reversal}.};
\node[box, below=of everett] (lineage) {\textbf{Diachronic inheritance}\\
Theorems~\ref{thm:identity-preserving-lineage} and \ref{thm:diachronic-practice}.};
\node[box, right=of lineage] (taxonomy) {\textbf{Interpretation exhaustion}\\
Theorem~\ref{thm:selector-taxonomy}.};
\draw[line] (substrate) -- (semantic);
\draw[line] (substrate) -- (everett);
\draw[line] (semantic) -- (confirmation);
\draw[line] (everett) -- (lineage);
\draw[line] (confirmation) -- (toy);
\draw[line] (confirmation) -- (taxonomy);
\draw[line] (lineage) -- (taxonomy);
\end{tikzpicture}
\caption{Dependency map for the hardened theorem ladder. The quotient theorem and the Everettian
non-selection theorem are the two main inputs to the confirmation and lineage theorems; the
interpretation taxonomy is downstream of both.}
\label{fig:dependency-map}
\end{figure}

\subsection{Denial-condition matrix}

\begin{longtable}{>{\raggedright\arraybackslash}p{0.28\textwidth} >{\raggedright\arraybackslash}p{0.40\textwidth} >{\raggedright\arraybackslash}p{0.20\textwidth}}
\caption{Denial-condition matrix for standard objections.}\\
\toprule
\textbf{Objection} & \textbf{To maintain it, one must deny} & \textbf{Where to look} \\
\midrule
\endfirsthead
\toprule
\textbf{Objection} & \textbf{To maintain it, one must deny} & \textbf{Where to look} \\
\midrule
\endhead
The framework is optional & The compressed substrate-necessity bridge or the standing conditional itself & Sec.~\ref{sec:targetclass} \\
\addlinespace
Different admissible assignments are just semantic packagings & The equivalence-collapse theorem or Corollary~\ref{cor:confirmational-divergence-quotient-relevant} & Sec.~\ref{sec:semanticcore}; App.~\ref{app:formal-quotient} \\
\addlinespace
Branch-local determinacy is enough & Lemma~\ref{thm:branchlocal-not-diachronic} & Sec.~\ref{sec:branching} \\
\addlinespace
Decision theory solves the problem & Theorem~\ref{thm:decision-theoretic-dependence} & Sec.~\ref{sec:diachronic_practice} \\
\addlinespace
Confirmation can be relativized without loss & Theorems~\ref{thm:theory-level-confirmation-criterion} and \ref{thm:fixation-commutation-necessity} & Sec.~\ref{sec:confirmation} \\
\addlinespace
Memory, anticipation, and agency are separate issues & Theorems~\ref{thm:identity-preserving-lineage} and \ref{thm:diachronic-inheritance} & Sec.~\ref{sec:diachronic_practice} \\
\addlinespace
There may be a fifth interpretive family & Theorem~\ref{thm:selector-exhaustion} & Sec.~\ref{sec:taxonomy} \\
\addlinespace
Instrumentalism or QBism refutes the theorem ladder & Proposition~\ref{prop:deflation-exit} & Sec.~\ref{sec:objections} \\
\bottomrule
\label{tab:attack-map}
\end{longtable}

\subsection{Compact reading path}

A short but faithful route through the manuscript is:
\begin{enumerate}
    \item Section~\ref{sec:targetclass}: fixed downstream substrate, target discourse, and standing
    conditional.
    \item Section~\ref{sec:semanticcore}: Definitions~\ref{def:admissible-reference-assignment}--\ref{def:mr-presentation},
    Theorem~\ref{thm:measurement-reference-quotient}, and
    Corollary~\ref{cor:confirmational-divergence-quotient-relevant}.
    \item Section~\ref{sec:branching}: Corollary~\ref{cor:everettian-nonselection} and
    Theorem~\ref{thm:everett_nonselection}.
    \item Section~\ref{sec:confirmation}: Theorems~\ref{thm:theory-level-confirmation-criterion},
    \ref{thm:fixation-commutation-necessity}, and
    \ref{thm:reference-stabilization-necessity}.
    \item Section~\ref{sec:toy}: Theorem~\ref{thm:toy-reversal}.
    \item Section~\ref{sec:diachronic_practice}: Theorems~\ref{thm:identity-preserving-lineage} and
    \ref{thm:diachronic-practice}.
    \item Section~\ref{sec:taxonomy}: Theorem~\ref{thm:selector-taxonomy}.
    \item Section~\ref{sec:objections}: only the objection subsections still thought live after the first
    seven steps.
\end{enumerate}

\subsection{Role of the appendix}

This appendix is procedural. The dependency map, denial-condition matrix, and reading path expose
the architecture of the argument without adding premises to it. Table~\ref{tab:strict-denial-table} in
Section~\ref{sec:objections} is the main-text compression of the same hostile-referee logic.
